\documentclass[10pt,a4paper]{amsart}
\usepackage{amsmath,amssymb,amsthm,graphicx}
% Retained preamble for the Stacks Project chapter snapshot in this directory.
% The chapter's own preamble.tex is a copy of the shared Stacks Project preamble; the block
% below is sliced MECHANICALLY, line for line, from it. Nothing was added, removed or
% redefined; the wrapper class, the xr-hyper/hyperref wiring and the externaldocument list
% are omitted because the assembler supplies the document class.
%
% lines taken from preamble.tex: 3, 5, 9, 11, 12, 156, 157, 158, ... (51 lines total)
%
% package coverage restored from the chapter's own preamble (2026-10-13): xy, xy-2cell

\usepackage[all]{xy}
\xyoption{2cell}
\UseAllTwocells

\newenvironment{reference}{\comment}{\endcomment}
\newenvironment{slogan}{\comment}{\endcomment}
\newenvironment{history}{\comment}{\endcomment}
\theoremstyle{plain}
\newtheorem{theorem}[subsection]{Theorem}
\newtheorem{proposition}[subsection]{Proposition}
\newtheorem{lemma}[subsection]{Lemma}
\theoremstyle{definition}
\newtheorem{definition}[subsection]{Definition}
\newtheorem{example}[subsection]{Example}
\newtheorem{exercise}[subsection]{Exercise}
\newtheorem{situation}[subsection]{Situation}
\theoremstyle{remark}
\newtheorem{remark}[subsection]{Remark}
\newtheorem{remarks}[subsection]{Remarks}
\numberwithin{equation}{subsection}
\def\lim{\mathop{\mathrm{lim}}\nolimits}
\def\colim{\mathop{\mathrm{colim}}\nolimits}
\def\Spec{\mathop{\mathrm{Spec}}}
\def\Hom{\mathop{\mathrm{Hom}}\nolimits}
\def\Ext{\mathop{\mathrm{Ext}}\nolimits}
\def\SheafHom{\mathop{\mathcal{H}\!\mathit{om}}\nolimits}
\def\SheafExt{\mathop{\mathcal{E}\!\mathit{xt}}\nolimits}
\def\Sch{\mathit{Sch}}
\def\Mor{\mathop{\mathrm{Mor}}\nolimits}
\def\Ob{\mathop{\mathrm{Ob}}\nolimits}
\def\Sh{\mathop{\mathit{Sh}}\nolimits}
\def\NL{\mathop{N\!L}\nolimits}
\def\CH{\mathop{\mathrm{CH}}\nolimits}
\def\proetale{{pro\text{-}\acute{e}tale}}
\def\etale{{\acute{e}tale}}
\def\QCoh{\mathit{QCoh}}
\def\Ker{\mathop{\mathrm{Ker}}}
\def\Im{\mathop{\mathrm{Im}}}
\def\Coker{\mathop{\mathrm{Coker}}}
\def\Coim{\mathop{\mathrm{Coim}}}
\DeclareMathSymbol{\boxtimes}{\mathbin}{AMSa}{"02}
\def\QCohstack{\mathcal{QC}\!\mathit{oh}}
\def\Cohstack{\mathcal{C}\!\mathit{oh}}
\def\Spacesstack{\mathcal{S}\!\mathit{paces}}
\def\Quotfunctor{\mathrm{Quot}}
\def\Hilbfunctor{\mathrm{Hilb}}
\def\Curvesstack{\mathcal{C}\!\mathit{urves}}
\def\Polarizedstack{\mathcal{P}\!\mathit{olarized}}
\def\Complexesstack{\mathcal{C}\!\mathit{omplexes}}
\def\Pic{\mathop{\mathrm{Pic}}\nolimits}
\def\Picardstack{\mathcal{P}\!\mathit{ic}}
\def\Picardfunctor{\mathrm{Pic}}
\def\Deformationcategory{\mathcal{D}\!\mathit{ef}}
\usepackage{cleveref}
\crefname{originallemma}{original lemma}{original lemmas}
\Crefname{originallemma}{Original lemma}{Original lemmas}
\makeatletter
\newlabel{algebra-lemma-henselization-different}{{04GV}{source}{}{sourcepointer-1}{}}
\newlabel{algebra-lemma-henselization-different@cref}{{[originallemma][0][]04GV}{[1][0][]source}{}{}{}}
\makeatother
\crefname{originallemma}{original lemma}{original lemmas}
\Crefname{originallemma}{Original lemma}{Original lemmas}
\makeatletter
\newlabel{algebra-lemma-strict-henselization-different}{{04GW}{source}{}{sourcepointer-2}{}}
\newlabel{algebra-lemma-strict-henselization-different@cref}{{[originallemma][0][]04GW}{[1][0][]source}{}{}{}}
\makeatother
\crefname{originallemma}{original lemma}{original lemmas}
\Crefname{originallemma}{Original lemma}{Original lemmas}
\makeatletter
\newlabel{lemma-bootstrap}{{08R1}{source}{}{sourcepointer-3}{}}
\newlabel{lemma-bootstrap@cref}{{[originallemma][0][]08R1}{[1][0][]source}{}{}{}}
\makeatother
\crefname{originallemma}{original lemma}{original lemmas}
\Crefname{originallemma}{Original lemma}{Original lemmas}
\makeatletter
\newlabel{lemma-colimit-cotangent-complex}{{08S9}{source}{}{sourcepointer-4}{}}
\newlabel{lemma-colimit-cotangent-complex@cref}{{[originallemma][0][]08S9}{[1][0][]source}{}{}{}}
\makeatother
\crefname{originallemma}{original lemma}{original lemmas}
\Crefname{originallemma}{Original lemma}{Original lemmas}
\makeatletter
\newlabel{lemma-derived-diagonal}{{08R0}{source}{}{sourcepointer-5}{}}
\newlabel{lemma-derived-diagonal@cref}{{[originallemma][0][]08R0}{[1][0][]source}{}{}{}}
\makeatother
\crefname{originaldefinition}{original definition}{original definitions}
\Crefname{originaldefinition}{Original definition}{Original definitions}
\makeatletter
\newlabel{more-algebra-definition-weakly-etale}{{092B}{source}{}{sourcepointer-6}{}}
\newlabel{more-algebra-definition-weakly-etale@cref}{{[originaldefinition][0][]092B}{[1][0][]source}{}{}{}}
\makeatother
\title{Nine cases in which the cotangent complex vanishes}
\author{The Stacks Project Authors}
\date{stacks-project:cotangent.tex:local-snapshot-20261012; GNU FDL 1.2 or later}
\begin{document}
\maketitle
\setcounter{subsection}{1}
\noindent\textit{Editorial note.} All delivered mathematical text is the author's own, copied from the pinned Stacks Project chapter file cotangent.tex; no statement or proof step was rewritten, reordered or removed.\par\medskip
\noindent\textit{Editorial note.} The statement's last item is the author's own placeholder entry 'add more here', kept verbatim inside the counted statement span as part of the author text.\par\medskip
\noindent\textit{Editorial note.} The six references of the counted spans are recorded as non-load-bearing navigation pointers; the argument that consumes them stays inside the counted proof span.\par\medskip
\noindent\textit{Editorial note.} Preamble: the retained preamble of this chapter (sources/stacks-cotangent/retained-preamble.tex, lines 1-63), containing the Xy-pic setup, the comment environments, the theorem declarations and the macros including \textbackslash{}NL.\par\medskip
\noindent\textit{Allowed standard prerequisites (editorial).}\begin{itemize}
\item The derived diagonal lemma for a surjective flat ring map. Lemma lemma-derived-diagonal, used at proof lines 1852-1855. Gives the vanishing over the tensor square in case (1).
\item The bootstrap lemma. Lemma lemma-bootstrap (1804-1809, produced as L44151), used at proof lines 1855-1857. Concludes the vanishing in case (1).
\item The cotangent complex commutes with filtered colimits. Lemma lemma-colimit-cotangent-complex (373-377, produced as L44165), used at proof lines 1858-1858. Gives case (6).
\item (Strict) henselizations as filtered colimits of etale ring extensions. Algebra, Lemmas algebra-lemma-henselization-different and algebra-lemma-strict-henselization-different, used at proof lines 1859-1862. Gives cases (7) and (8).
\end{itemize}\medskip
\ifdefined\hypertarget\hypertarget{sourcepointer-1}{}\fi
\noindent\textit{Source reference.} Original lemma 04GV: focused-seeds/AGP/sources/stacks-snapshot-20261011/tags.txt, source lines 5846--5846. Navigation pointer into the Stacks Project; the standard fact it supplies is applied inside the counted proof span and the argument resting on it stays there.\par\medskip
\ifdefined\hypertarget\hypertarget{sourcepointer-2}{}\fi
\noindent\textit{Source reference.} Original lemma 04GW: focused-seeds/AGP/sources/stacks-snapshot-20261011/tags.txt, source lines 5847--5847. Navigation pointer into the Stacks Project; the standard fact it supplies is applied inside the counted proof span and the argument resting on it stays there.\par\medskip
\ifdefined\hypertarget\hypertarget{sourcepointer-3}{}\fi
\noindent\textit{Source reference.} Original lemma 08R1: focused-seeds/AGP/sources/stacks-snapshot-20261011/tags.txt, source lines 11079--11079. Navigation pointer into the Stacks Project; the standard fact it supplies is applied inside the counted proof span and the argument resting on it stays there.\par\medskip
\ifdefined\hypertarget\hypertarget{sourcepointer-4}{}\fi
\noindent\textit{Source reference.} Original lemma 08S9: focused-seeds/AGP/sources/stacks-snapshot-20261011/tags.txt, source lines 11122--11122. Navigation pointer into the Stacks Project; the standard fact it supplies is applied inside the counted proof span and the argument resting on it stays there.\par\medskip
\ifdefined\hypertarget\hypertarget{sourcepointer-5}{}\fi
\noindent\textit{Source reference.} Original lemma 08R0: focused-seeds/AGP/sources/stacks-snapshot-20261011/tags.txt, source lines 11078--11078. Navigation pointer into the Stacks Project; the standard fact it supplies is applied inside the counted proof span and the argument resting on it stays there.\par\medskip
\ifdefined\hypertarget\hypertarget{sourcepointer-6}{}\fi
\noindent\textit{Source reference.} Original definition 092B: focused-seeds/AGP/sources/stacks-snapshot-20261011/tags.txt, source lines 11474--11474. Navigation pointer into the Stacks Project; the standard fact it supplies is applied inside the counted proof span and the argument resting on it stays there.\par\medskip
\begin{lemma}
\label{lemma-when-zero}
The cotangent complex $L_{B/A}$ is zero in each of the following cases:
\begin{enumerate}
\item $A \to B$ and $B \otimes_A B \to B$ are flat, i.e., $A \to B$
is weakly \'etale
(More on Algebra, Definition \ref{more-algebra-definition-weakly-etale}),
\item $A \to B$ is a flat epimorphism of rings,
\item $B = S^{-1}A$ for some multiplicative subset $S \subset A$,
\item $A \to B$ is unramified and flat,
\item $A \to B$ is \'etale,
\item $A \to B$ is a filtered colimit of ring maps for which
the cotangent complex vanishes,
\item $B$ is a henselization of a local ring of $A$,
\item $B$ is a strict henselization of a local ring of $A$, and
\item add more here.
\end{enumerate}
\end{lemma}

\begin{proof}
In case (1) we may apply
Lemma \ref{lemma-derived-diagonal}
to the surjective flat ring map $B \otimes_A B \to B$
to conclude that $L_{B/B \otimes_A B} = 0$ and then we use
Lemma \ref{lemma-bootstrap}
to conclude. The cases (2) -- (5) are each special cases of (1).
Part (6) follows from Lemma \ref{lemma-colimit-cotangent-complex}.
Parts (7) and (8) follows from the fact that (strict) henselizations
are filtered colimits of \'etale ring extensions of $A$, see
Algebra, Lemmas \ref{algebra-lemma-henselization-different} and
\ref{algebra-lemma-strict-henselization-different}.
\end{proof}

\end{document}
